% Teil II: durchgerechneter Weg, Teile a--c.

\subsection*{Gegeben und gesucht}
Tröte der Länge $l$, ein Ende geschlossen, ein Ende offen, gef"ullt mit
Stickstoff N$_2$. Normalbedingungen $T=\qty{0}{\celsius}=273.15$~K,
$p=1$~atm. $R=\qty{8.314}{J/(mol\,K)}$, $\kappa_{\mathrm{N}_2}=7/5$,
$\kappa_{\mathrm{He}}=5/3$, $M_{\mathrm{N}_2}=\qty{28}{g/mol}$,
$M_{\mathrm{He}}=\qty{4}{g/mol}$. Gesucht: (a) Skizze Grundschwingung +
erste Oberschwingung; (b) Frequenz der Grundschwingung; (c) "Anderung bei
Ersatz durch Helium.

\subsection*{Ansatz und Physik dahinter}
Einseitig geschlossenes Rohr $\Rightarrow$ Grundwelle $l=\lambda/4$:
$f_1=c/4l$. Schallgeschwindigkeit im idealen Gas $c=\sqrt{\kappa R T/M}$.
Die Länge $l$ ist nicht gegeben --- das Ergebnis bleibt symbolisch in $l$,
aber $c$ l"asst sich mit den gegebenen Konstanten ausrechnen. Für Teil (c)
ist nur das \emph{Verh"altnis} der beiden Schallgeschwindigkeiten n"otig,
da die Geometrie ($l$) unver"andert bleibt: $f_{\mathrm{He}}/f_{\mathrm{N}_2}
=c_{\mathrm{He}}/c_{\mathrm{N}_2}$.

\subsection*{Rechnung}
\textbf{a) Skizzen.} Grundschwingung: Knoten am geschlossenen Ende, Bauch am
offenen Ende, genau eine Viertelwelle ($\lambda_1=4l$). Erste Oberschwingung:
drei Viertelwellen ($\lambda=4l/3$, $f=3f_1$); Knoten am geschlossenen Ende,
ein Knoten bei $l/3$, Bauch bei $2l/3$ und am offenen Ende.

\textbf{b) Grundfrequenz (N$_2$).} Zuerst $c$ (mit M in \unit{kg/mol}):
\begin{equation}
  c_{\mathrm{N}_2}
  =\sqrt{\frac{\kappa R\,T}{M}}
  =\sqrt{\frac{7/5\cdot8.314\cdot273.15}{0.028}}
  =\sqrt{\frac{3180.6}{0.028}}
  \approx\boxed{\;\qty{337}{m/s}\;}.
\end{equation}
Damit
\begin{equation}
  \boxed{\;f_1=\frac{c_{\mathrm{N}_2}}{4l}
  =\frac{337}{4\,l}\,\mathrm{m/s}\;\bigl(l\ \text{in } m\bigr)
  \approx\frac{84.3}{l}\,\mathrm{Hz}\;}.
\end{equation}
Also z. B. f"ur $l=\qty{1}{m}$: $f_1\approx\qty{84}{Hz}$.

\textbf{c) Ersatz durch Helium.} Nur das Verh"altnis der $c$'s z"ahlt:
\begin{equation}
  \frac{c_{\mathrm{He}}}{c_{\mathrm{N}_2}}
  =\sqrt{\frac{\kappa_{\mathrm{He}}}{\kappa_{\mathrm{N}_2}}
  \cdot\frac{M_{\mathrm{N}_2}}{M_{\mathrm{He}}}}
  =\sqrt{\frac{5/3}{7/5}\cdot\frac{28}{4}}
  =\sqrt{\frac{25}{3}}
  \approx\boxed{\;2.89\;}.
\end{equation}
Zahlenprobe ($c_{\mathrm{He}}=\sqrt{5/3\cdot8.314\cdot273.15/0.004}
\approx\qty{973}{m/s}$, geteilt durch 337 $\approx2.89$). Die
Grundfrequenz wird also um den Faktor $2.89$ \emph{h"oher}:
\begin{equation}
  f_{1,\mathrm{He}}\approx 2.89\,f_{1,\mathrm{N}_2}\approx\frac{243}{l}
  \,\mathrm{Hz}\qquad(l\ \text{in } m).
\end{equation}

\begin{ergebnis}
(a) Grundschwingung: Halbwelle $\lambda/4$ (Knoten am einen, Bauch am
offenen Ende); \checkmark
erste Oberschwingung: $\lambda=4l/3$, mit zus"atzlichem Knoten.
(b) $c_{\mathrm{N}_2}=\qty{337}{m/s}$,
$f_1=c_{\mathrm{N}_2}/4l=84.3/l$ Hz.
(c) Helium: $c_{\mathrm{He}}=973$~m/s, Faktor $2.89$,
$f_{1,\mathrm{He}}\approx243/l$ Hz --- h"oherer Ton.
\end{ergebnis}

\begin{kontrolle}
\textbf{Einheiten:} $\sqrt{\kappa R T/M}$:
$\sqrt{(\mathrm{kg\,m^2/s^2})/(\mathrm{kg})}=\mathrm{m/s}$. \checkmark

\textbf{Größenordnung:} $c_{\mathrm{N}_2}=337$~m/s ist nur wenig kleiner als
Schall in Luft ($\approx343$~m/s), weil N$_2$ und die effektive molare Masse
der Luft fast gleich sind. Das passt.\ \checkmark

\textbf{Konsistenz von Teil (b,c):} da $f_1\propto c\propto1/\sqrt M$, und Helium
sehr leicht ist, ist eine "Anderung um Faktor $\approx2.9$ plausibel.
\checkmark

\textbf{Grenzfall:} bei $T\to0$ ginge $c\to0$ (alle Molek"ule st"anden
still)---im klassischen Bild sinnvoll; die Formel verliert erst weit unter
Konsensbedingungen ihre G"ultigkeit.
\end{kontrolle}

\begin{merksatz}
Einseitig geschlossenes Rohr (Tröte): Grundwelle $l=\lambda/4$, $f_1=c/4l$,
Obertöne nur ungerade Vielfache. Im idealen Gas
$c=\sqrt{\kappa R T/M}$; wegen $\kappa\propto\sqrt{\kappa/M}$ wird ein
leichtes einatomiges Gas (Helium) deutlich "schneller'' als schweres,
zweiatomiges Stickstoff, die Tröte klingt h"oher.
\end{merksatz}
